\documentclass[12pt]{article}
\usepackage[T1]{fontenc}
\usepackage[utf8]{inputenc}
\usepackage{lmodern}
\usepackage{textcomp}
\usepackage{enumitem}
\usepackage{geometry}
\geometry{margin=1in}
\begin{document}
\begin{center}
Answer Key - Business Administration - Graduate Level Assessment 18 \
\small (Answers are ordered exactly as in the question paper.)
\end{center}

\section*{Multiple Choice}
\begin{enumerate}[label=\arabic*.]
\item \textbf{Answer:} A
\\ \textit{Options:} A) 1800.0000; B) 2500.0000; C) 2200.5678; D) 2300.4567; E) 2600.8910
\\ \textit{Note:} The present value of the three end-of-year payments of $1000 each, starting from the end of the fifth year at a 7% rate of return, is calculated using the present value formula, which discounts future cash flows back to the present, resulting in a total present value of $1800.00.
\item \textbf{Answer:} E
\\ \textit{Options:} A) \$212.36; B) \$220.50; C) \$210.00; D) \$222.22; E) \$216.49
\\ \textit{Note:} The correct answer of $216.49 is derived from applying the formula for compound interest, which accounts for the principal amount of $200, an interest rate of 4\% compounded semiannually over 2 years, resulting in a future value calculation of \$200 * (1 + 0.04/$2)^(2$*2).
\item \textbf{Answer:} A
\\ \textit{Options:} A) \$325.50; B) \$650; C) \$585; D) \$405; E) \$975
\\ \textit{Note:} The correct answer is $325.50 because George earns a commission of $22.50 per vacuum cleaner (30\% of $75), and with 13 sales, his total earnings amount to $22.50 multiplied by 13, which equals $292.50, plus the base salary of $33, resulting in \$325.50.
\item \textbf{Answer:} A
\\ \textit{Options:} A) Individual skills; B) Cultural influences; C) Economic conditions; D) Political situation; E) Personality traits
\\ \textit{Note:} The situational theory of leadership emphasizes the importance of adapting one's leadership style based on the individual skills and readiness of team members to effectively guide and support them in various contexts.
\item \textbf{Answer:} A
\\ \textit{Options:} A) Operant conditioning.; B) Social learning.; C) Insight learning.; D) Observational learning.; E) Cognitive learning.
\\ \textit{Note:} Providing free samples of perfumes in magazines exemplifies operant conditioning because it uses positive reinforcement by rewarding potential customers with a free product to encourage future purchases.
\end{enumerate}

\section*{True / False}
\begin{enumerate}[label=\arabic*.]
\setcounter{enumi}{5}
\item \textbf{Answer:} False
\\ \textit{Options:} A) True; B) False
\\ \textit{Note:} The statement is false because investing $1,000 to receive $150 annually for 10 years results in a return of 15\% per year, not 10\%, as the total cash inflow over the period exceeds the initial investment when calculated using the present value of annuities.
\item \textbf{Answer:} False
\\ \textit{Options:} A) True; B) False
\\ \textit{Note:} The statement is false because the book value per share is calculated by dividing total equity by the number of outstanding shares, and without knowing the total equity or the number of shares, we cannot determine the book value per share accurately.
\item \textbf{Answer:} False
\\ \textit{Options:} A) True; B) False
\\ \textit{Note:} The statement is false because egoism and utilitarianism are both consequentialist theories that evaluate the morality of actions based on their outcomes rather than adhering to fixed moral duties.
\item \textbf{Answer:} True
\\ \textit{Options:} A) True; B) False
\\ \textit{Note:} The statement is true because the terms (6/10), (n/30) indicate a 6\% discount available if paid within 10 days, and when calculated for the entire year, this translates to an effective annual percentage rate of approximately 24.49\%.
\item \textbf{Answer:} True
\\ \textit{Options:} A) True; B) False
\\ \textit{Note:} The statement is true because it accurately reflects that the total payment made by Mr. Scaccioborrows encompasses both the principal amount borrowed and the interest accrued over the repayment period, which is a standard practice in loan agreements.
\end{enumerate}

\section*{Long Form Response}
\begin{enumerate}[label=\arabic*.]
\setcounter{enumi}{10}
\item \textbf{Answer:} The theoretical forward price of gold can be calculated by considering the current price, storage costs, and the opportunity cost represented by the interest rate. Given a current price of $412 per ounce, an interest rate of 9% compounded quarterly for 9 months, and assuming storage costs are negligible for this calculation, the forward price can be derived using the formula: Forward Price = Spot Price * (1 + r/n)^(nt), where r is the interest rate, n is the number of compounding periods per year, and t is the time in years. Substituting the values, we find that the forward price is approximately $442.02 per ounce. This price reflects the cost of carrying the gold over the specified period, accounting for the time value of money. Thus, the theoretical forward price is influenced by both the current market conditions and the cost of financing the purchase of gold.
\\ \textit{Note:} correct because it effectively applies the formula for calculating the theoretical forward price of gold by incorporating the current spot price, the compounded interest rate over the specified period, and implicitly assumes negligible storage costs, which are critical in determining the future price in a forward contract.
\item \textbf{Answer:} Sequential discounts can significantly impact the final pricing of a product, as they are applied successively rather than cumulatively. For a Speedway Racing Set priced at $84, the first discount of 20% reduces the price to $67.20. Applying the subsequent 10\% discount to this new amount further lowers the price to \$60.48. This approach demonstrates that the final price is not merely the sum of the two discounts but rather reflects a compounding effect where each discount is applied to the remaining balance. Understanding this pricing strategy is crucial for businesses to effectively manage profit margins and consumer perceptions of value.
\\ \textit{Note:} correct because it illustrates how sequential discounts reduce the price stepwise, showing that each discount is applied to the already reduced amount rather than cumulatively, resulting in a net price that reflects the compounding effect of the discounts.
\end{enumerate}

\vfill
\noindent\textit{End of Answer Key}
\end{document}